\CWHeader{Laboratory works No. 2 and No. 3: PATRICIA and Quality Research}

\begin{table}[!ht]
\centering
\begin{tabular}{|l|l|}
\hline
Time limit & 15 seconds \\
\hline
Memory limit & 512 MB \\
\hline
Input & standard input or \texttt{input.txt} \\
\hline
Output & standard output or \texttt{output.txt} \\
\hline
\end{tabular}
\end{table}

\section*{Laboratory work No. 2}

\CWProblem{
You need to create a program library implementing the specified data structure, on the basis of
which to develop a dictionary program. In the dictionary, each key, which is a case-insensitive
sequence of letters of the English alphabet of no more than 256 characters, is associated with a
number from 0 to $2^{64} - 1$. Different words can be assigned the same number.
}

The program must process lines of the input file until its end. Each line can have the following
format:

\begin{itemize}
    \item \texttt{+ word 34}~--- add the word \texttt{word} with number 34 to the dictionary. The
    program must output \texttt{OK} if the operation was successful, \texttt{Exist} if the word is
    already in the dictionary.
    \item \texttt{- word}~--- remove the word \texttt{word} from the dictionary. The program must
    output \texttt{OK} if the word existed and was removed, \texttt{NoSuchWord} if the word was not
    found in the dictionary.
    \item \texttt{word}~--- find the word \texttt{word} in the dictionary. The program must output
    \texttt{OK: 34} if the word was found. If the word was not found in the dictionary, output
    \texttt{NoSuchWord}.
    \item \texttt{! Save /path/to/file}~--- save the dictionary to a file in binary compact
    representation, at the path specified by the command parameter. If successful, the program
    must output \texttt{OK}; if the operation fails, output a description of the error.
    \item \texttt{! Load /path/to/file}~--- load the dictionary from a file. It is assumed that the
    file was previously prepared using the \texttt{Save} command. If successful, the program must
    output \texttt{OK} and the loaded dictionary replaces the current one; if unsuccessful,
    diagnostics must be output and the working dictionary must remain unchanged.
\end{itemize}

For all operations, if a system error occurs, the program must output a line starting with
\texttt{ERROR:} and describing the error in English.

The difference between variants lies only in the data structures used.

{\bfseries Data structure variant:} PATRICIA.

\section*{Laboratory work No. 3: A. Software Quality Research}

\CWProblem{
For the dictionary implementation from the previous laboratory work, it is necessary to
investigate execution speed and RAM consumption. If errors or obvious shortcomings are found, they
must be fixed.
}

The result of the work is a report containing:

\begin{itemize}
    \item a work diary describing what was done and when, which tools were used, and what results
    were achieved at each step of the laboratory work;
    \item conclusions about the shortcomings found;
    \item comparison of the fixed program with the previous version;
    \item general conclusions about the completed laboratory work and the experience gained.
\end{itemize}

The minimum set of tools used must include the \texttt{gprof} utility and the \texttt{dmalloc}
library. However, they may be replaced with any similar or more advanced tools, for example
Valgrind or Shark, and additional tools may also be used, for example \texttt{gcov}.

\subsection*{Additional information}

\subsubsection*{gprof}

The \texttt{gprof} utility should be studied independently using the operational documentation of
the Unix operating system.

\subsubsection*{dmalloc}

Below is a short description of how to use it for a small test program.

Suppose a small program is written with two dynamic memory usage errors that cannot be checked at
compile time and can only be detected at run time:

\begin{verbatim}
#include <stdio.h>
#include <stdlib.h>

#ifdef DMALLOC
#include "dmalloc.h"
#endif

void *foo() {
    char *p = malloc(1023);
    p[1023] = '\0';
    return p;
}

int main() {
    foo();
    return 0;
}
\end{verbatim}

The program has two errors:

\begin{enumerate}
    \item In line 10, the program accesses the 1024th element of the dynamic array \texttt{p}, even
    though only 1023 elements were allocated. When the program is executed, this error will most
    likely not be detected, because the actual allocated memory size is usually aligned to the page
    size and therefore will be larger than 1023 elements.
    \item The memory allocated in line 9 is never freed.
\end{enumerate}

If the program is compiled and run, most likely no problems will occur either during compilation
or during execution:

\begin{verbatim}
gcc -Wall dmalloc-test.c -o dtest
./dtest
\end{verbatim}

To use the \texttt{dmalloc} library, the special header file \texttt{dmalloc.h} must be included in
the program. It replaces calls to memory allocation functions with its own versions. The
\texttt{DMALLOC} macro must also be defined to enable the main functionality of the library. In
addition, the \texttt{DMALLOC\_FUNC\_CHECK} macro can be defined to enable additional checking of
parameters passed to various functions from the C standard library. Finally, the \texttt{dmalloc}
library must be linked during compilation.

Thus, compilation can be performed as follows:

\begin{verbatim}
gcc -Wall -DDMALLOC -DDMALLOC_FUNC_CHECK \
    -I/usr/local/include -L/usr/local/lib dmalloc-test.c \
    -o dtest -ldmalloc
\end{verbatim}

Before running the program under investigation, the \texttt{DMALLOC\_OPTIONS} environment variable
must be defined. It passes information about how the library should work. The full list of options
and their meaning should be found in the library documentation. To correctly set the environment
in the \texttt{bash} command interpreter, the following command can be used:

\begin{verbatim}
eval `command dmalloc -b -l logfile -i 100 low`
\end{verbatim}

In this case, the check report will be written to the file \texttt{logfile}.

Now, if the program is run, it will terminate abnormally:

\begin{verbatim}
./dtest
\end{verbatim}

\begin{verbatim}
debug-malloc library: dumping program, fatal error
   Error: failed OVER picket-fence magic-number check (err 27)
Abort trap: 6 (core dumped)
\end{verbatim}

Thus, the \texttt{dmalloc} library detected an out-of-bounds access. The abnormal program
termination also made it possible to create a core dump, which can be examined using the
\texttt{gdb} debugger.

If the source program is modified by fixing the out-of-bounds access, for example by replacing
index 1023 with 1022, the program will run normally. The report will look approximately as
follows:

{\scriptsize
\begin{alltt}
1192785155: 1: Dmalloc version '5.5.2' from 'http://dmalloc.com/'
1192785155: 1: flags = 0x4e48503, logfile 'logfile'
1192785155: 1: interval = 100, addr = 0, seen # = 0, limit = 0
1192785155: 1: starting time = 1192785155
1192785155: 1: process pid = 23403
1192785155: 1: Dumping Chunk Statistics:
1192785155: 1: basic-block 8192 bytes, alignment 8 bytes
1192785155: 1: heap address range: 0x16005a000 to 0x160060000, 24576 bytes
1192785155: 1:
1192785155: 1:
1192785155: 1:
1192785155: 1: heap checked 1
1192785155: 1: alloc calls: malloc 1, calloc 0, realloc 0, free 0
1192785155: 1: alloc calls: recalloc 0, memalign 0, valloc 0
1192785155: 1: alloc calls: new 0, delete 0
 user blocks: 1 blocks, 7167 bytes (29\%)
admin blocks: 2 blocks, 16384 bytes (66\%)
total blocks: 3 blocks, 24576 bytes
1192785155: 1:   current memory in use: 1023 bytes (1 pnts)
1192785155: 1:  total memory allocated: 1023 bytes (1 pnts)
1192785155: 1:  max in use at one time: 1023 bytes (1 pnts)
1192785155: 1: max alloced with 1 call: 1023 bytes
1192785155: 1: max unused memory space: 1025 bytes (50\%)
1192785155: 1: top 10 allocations:
1192785155: 1: total-size count in-use-size count source
1192785155: 1: 1023 1 1023 1 dmalloc-test.c:9
1192785155: 1: 1023 1 1023 1 Total of 1
1192785155: 1: Dumping Not-Freed Pointers Changed Since Start:
1192785155: 1: not freed: '0x16005b808|s1' (1023 bytes) from 'dmalloc-test.c:9'
1192785155: 1: total-size count source
1192785155: 1: 1023 1 dmalloc-test.c:9
1192785155: 1: 1023 1 Total of 1
1192785155: 1: ending time = 1192785155, elapsed since start = 0:00:00
\end{alltt}
}

At the end of the report there is a \texttt{Dumping Not-Freed Pointers} section, where it can be
seen that one memory block of 1023 bytes, allocated in line 9, was not freed before the program
finished. Thus, memory leaks in programs can be tracked.

Before performing the laboratory work, it is strongly recommended to independently study the
documentation for the \texttt{dmalloc} library. Using additional options not described above will
improve the impression of the completed work.

In this laboratory work, \texttt{dmalloc} is replaced by more advanced tools: Valgrind Memcheck and
Valgrind Massif. The research also uses \texttt{gprof}, a test generator, and a benchmark that
compares PATRICIA with \texttt{std::map}.

\subsection*{Input format}

Each non-empty line of the input contains one of the commands described above.

\subsection*{Output format}

For each input line, the program outputs one response line as specified above.

\subsection*{Example}

Input:

\begin{alltt}
+ a 1
+ A 2
+ \seqsplit{aaaaaaaaaaaaaaaaaaaaaaaaaaaaaaaaaaaaaaaaaaaaaaaaaaaaaaaaaaaaaaaaaaaaaaaaaaaaaaaaaaaaaaaaaaaaaaaaaaaaaaaaaaaaaaaaaaaaaaaaaaaaaaaaaaaaaaaaaaaaaaaaaaaaaaaaaaaaaaaaaaaaaaaaaaaaaaaaaaaaaaaaaaaaaaaaaaaaaaaaaaaaaaaaaaaaaaaaaaaaaaaaaaaaaaaaaaaaaaaaaaaaaaaaaaaaaaaa} 18446744073709551615
\seqsplit{aaaaaaaaaaaaaaaaaaaaaaaaaaaaaaaaaaaaaaaaaaaaaaaaaaaaaaaaaaaaaaaaaaaaaaaaaaaaaaaaaaaaaaaaaaaaaaaaaaaaaaaaaaaaaaaaaaaaaaaaaaaaaaaaaaaaaaaaaaaaaaaaaaaaaaaaaaaaaaaaaaaaaaaaaaaaaaaaaaaaaaaaaaaaaaaaaaaaaaaaaaaaaaaaaaaaaaaaaaaaaaaaaaaaaaaaaaaaaaaaaaaaaaaaaaaaaaaa}
A
- A
a
\end{alltt}

Output:

\begin{alltt}
OK
Exist
OK
OK: 18446744073709551615
OK: 1
OK
NoSuchWord
\end{alltt}

\pagebreak
